\section{Threats to Validity}

\paragraph{Construct validity.}
WER does not directly measure correction burden, command errors, editing disruption, or the
subjective cost of catastrophic inserted text. LibriSpeech is read audiobook speech rather than
spontaneous desktop dictation. Hosted-runner latency should not be treated as a portable measure
of interactive responsiveness. Synthetic TTS meeting fixtures are useful regression tests but
are not evidence of real-meeting diarization quality.

\paragraph{Internal validity.}
Several mechanism probes were designed after anomalies were observed, so they are explanatory
follow-ups rather than preregistered primary endpoints. Dependency versions differ across parts
of the campaign. Host load may interact with decode timing and fallback behavior, although the
available repeated runs do not establish host load as the cause of the \texttt{large-v3}
instability. Some experiments, particularly streaming, use small sample sizes.

\paragraph{External validity.}
The principal ASR evaluation is English-centric. The campaign does not yet provide a
representative spontaneous-speech study across accents, microphones, rooms, and noise levels.
The AMI headset mix is cleaner than a typical far-field laptop/table microphone. The study is
CPU-focused and does not characterize GPU execution or battery/energy cost.

\paragraph{Statistical validity.}
The campaign contains exploratory comparisons in addition to planned benchmark cells. We
therefore avoid inferring superiority from overlapping point estimates, use paired analyses when
conditions share utterances, and keep multiplicity corrections explicit in the onset analysis.
The diarization paired comparison contains only 16 meetings, limiting power to resolve modest
DER effects. The \texttt{large-v3} aggregate conditioning effect is concentrated in a few
utterances, so the paper reports both the mean result and that concentration.

\paragraph{Human factors.}
No controlled typing-versus-\YazSes{} study has yet measured final corrected entry rate,
correction time, or workload. The present work therefore makes no productivity claim relative to
typing. Likewise, technical dysfluency-filter results are not treated as evidence of benefit for
people who stutter or have dysarthria without a participant study.
